\markdownRendererDocumentBegin
\markdownRendererWarning{The `hybrid` option has been soft-deprecated.}{The `hybrid` option has been soft-deprecated.}{Consider using one of the following better options for mixing TeX and markdown: `contentBlocks`, `rawAttribute`, `texComments`, `texMathDollars`, `texMathSingleBackslash`, and `texMathDoubleBackslash`. For more information, see the user manual at <https://witiko.github.io/markdown/>.}{Consider using one of the following better options for mixing TeX and markdown: `contentBlocks`, `rawAttribute`, `texComments`, `texMathDollars`, `texMathSingleBackslash`, and `texMathDoubleBackslash`. For more information, see the user manual at <https://witiko.github.io/markdown/>.}\markdownRendererSectionBegin
\markdownRendererHeadingOne{Estructuras de Datos}\markdownRendererInterblockSeparator
{}\markdownRendererSectionBegin
\markdownRendererHeadingTwo{Arreglos Dinámicos}\markdownRendererInterblockSeparator
{}Un arreglo dinámico es un arreglo cuyo tamaño puede cambiar durante la ejecución del programa.\markdownRendererParagraphSeparator
{}La biblioteca estándar de C++ (STL) proporciona varias estructuras de arreglos dinámicos, siendo la más útil el vector.\markdownRendererInterblockSeparator
{}\markdownRendererSectionBegin
\markdownRendererHeadingThree{Vectores}\markdownRendererInterblockSeparator
{}Un vector permite agregar y eliminar elementos eficientemente al final de la estructura.\markdownRendererParagraphSeparator
{}Inicialización de un vector vacío\markdownRendererInterblockSeparator
{}\markdownRendererInputFencedCode{./_markdown_main/d09f662249e9ddca9aa585cdafb6c662.verbatim}{cpp}{cpp}\markdownRendererInterblockSeparator
{}Inicialización de un vector con lista de elementos\markdownRendererInterblockSeparator
{}\markdownRendererInputFencedCode{./_markdown_main/7628db821fa22993320e5d564df7086d.verbatim}{cpp}{cpp}\markdownRendererInterblockSeparator
{}Inicialización de un vector determinando el tamaño y valores iniciales\markdownRendererInterblockSeparator
{}\markdownRendererInputFencedCode{./_markdown_main/7fc1591f28a83235dee944219665a583.verbatim}{cpp}{cpp}\markdownRendererInterblockSeparator
{}Los elementos pueden ser accedidos de manera ordinaria\markdownRendererInterblockSeparator
{}\markdownRendererInputFencedCode{./_markdown_main/3ae209fe64be3f70babb3893e74e1a3e.verbatim}{cpp}{cpp}\markdownRendererInterblockSeparator
{}$\clearpage$\markdownRendererParagraphSeparator
{}La función \markdownRendererCodeSpan{size()} retorna el número de elementos en el vector\markdownRendererInterblockSeparator
{}\markdownRendererInputFencedCode{./_markdown_main/cbd5a992ac73f65f0dfef8b345eb55f8.verbatim}{cpp}{cpp}\markdownRendererInterblockSeparator
{}Manera más corta de iterar a través de un vector\markdownRendererInterblockSeparator
{}\markdownRendererInputFencedCode{./_markdown_main/85d773f3414d304e3335cae1c5583631.verbatim}{cpp}{cpp}\markdownRendererInterblockSeparator
{}La función \markdownRendererCodeSpan{back()} retorna el útlimo elemento de un vector, y la función \markdownRendererCodeSpan{pop\markdownRendererUnderscore{}back()} elimina el útlimo elemento\markdownRendererInterblockSeparator
{}\markdownRendererInputFencedCode{./_markdown_main/3b24ec4e30e341805a9fc0c22dac9f7c.verbatim}{cpp}{cpp}\markdownRendererInterblockSeparator
{}Los vectores se implementan de manera que las operaciones \markdownRendererCodeSpan{push\markdownRendererUnderscore{}back()} y \markdownRendererCodeSpan{pop\markdownRendererUnderscore{}back()} trabajen en O(1) tiempo medio.\markdownRendererParagraphSeparator
{}En la práctica, el uso de un vector es casi tan rápido como el uso una arreglo ordinario.\markdownRendererParagraphSeparator
{}
\markdownRendererSectionEnd \markdownRendererSectionBegin
\markdownRendererHeadingThree{Iteradores y rangos}\markdownRendererInterblockSeparator
{}Los iteradores son objetos que permiten recorrer los elementos de una estructura de datos de forma genérica.\markdownRendererSoftLineBreak
{}Funcionan como punteros que apuntan a elementos dentro de contenedores de la STL, como vector, set o map.\markdownRendererInterblockSeparator
{}\markdownRendererSectionBegin
\markdownRendererHeadingFour{Iterador}\markdownRendererInterblockSeparator
{}Es una variable que apunta a un elemento de una estructura de datos.\markdownRendererInterblockSeparator
{}\markdownRendererUlBegin
\markdownRendererUlItem El iterador \markdownRendererCodeSpan{begin()} apunta al primer elemento de una estructura de datos.\markdownRendererUlItemEnd 
\markdownRendererUlItem El iterador \markdownRendererCodeSpan{end()} apunta a la posición después del último elemento\markdownRendererUlItemEnd 
\markdownRendererUlEnd \markdownRendererInterblockSeparator
{}Los iteradores permiten recorrer estructuras sin conocer su implementación interna.\markdownRendererParagraphSeparator
{}El elemento al que apunta un iterador se puede acceder con la sintaxis *.\markdownRendererInterblockSeparator
{}\markdownRendererInputFencedCode{./_markdown_main/b51c83aca351d2f16c87af5ebbf70cd5.verbatim}{cpp}{cpp}\markdownRendererInterblockSeparator
{}
\markdownRendererSectionEnd \markdownRendererSectionBegin
\markdownRendererHeadingFour{Rango}\markdownRendererInterblockSeparator
{}Es una secuencia de elementos consecutivos en una estructura de datos.\markdownRendererParagraphSeparator
{}Los iteradores begin() y end() definen un rango que contiene todos los elementos de una estructura de datos.\markdownRendererParagraphSeparator
{}El código siguiente ordena primero un vector, luego invierte el orden de sus elementos, y por último, mezcla sus elementos.\markdownRendererInterblockSeparator
{}\markdownRendererInputFencedCode{./_markdown_main/92ab18c6e648dc1be6c5f863a893f5f8.verbatim}{cpp}{cpp}\markdownRendererInterblockSeparator
{}Un ejemplo más útil,\markdownRendererCodeSpan{lower\markdownRendererUnderscore{}bound} da un iterador al primer elemento en un rango ordenado cuyo valor es al menos x, y \markdownRendererCodeSpan{upper\markdownRendererUnderscore{}bound} da una\markdownRendererSoftLineBreak
{}iterador al primer elemento cuyo valor es mayor que x\markdownRendererInterblockSeparator
{}\markdownRendererInputFencedCode{./_markdown_main/aece67ad7a78d4a22f42e943f01c86c4.verbatim}{cpp}{cpp}\markdownRendererInterblockSeparator
{}El rango dado debe estar ordenado.\markdownRendererInterblockSeparator
{}
\markdownRendererSectionEnd 
\markdownRendererSectionEnd \markdownRendererSectionBegin
\markdownRendererHeadingThree{Deque}\markdownRendererInterblockSeparator
{}Un deque es un arreglo dinámico que puede manipularse eficientemente en ambos extremos.\markdownRendererParagraphSeparator
{}Los deques deben utilizarse si hay necesidad de manipular ambos extremos del arreglo.\markdownRendererInterblockSeparator
{}\markdownRendererInputFencedCode{./_markdown_main/4fecccaa0bc9c58e9098a8326c8882e4.verbatim}{cpp}{cpp}\markdownRendererParagraphSeparator
{}
\markdownRendererSectionEnd \markdownRendererSectionBegin
\markdownRendererHeadingThree{Stack}\markdownRendererInterblockSeparator
{}Una pila tiene las funciones push y pop para insertar y quitar elementos al final de la estructura, y la función top que recupera el último elemento\markdownRendererInterblockSeparator
{}\markdownRendererInputFencedCode{./_markdown_main/c5898f03e06b0fcf005011ed39eba119.verbatim}{cpp}{cpp}\markdownRendererInterblockSeparator
{}
\markdownRendererSectionEnd \markdownRendererSectionBegin
\markdownRendererHeadingThree{Queue}\markdownRendererInterblockSeparator
{}En una cola, los elementos se insertan al final de la estructura y se eliminan desde la parte delantera de la estructura. Ambas funciones front y back permiten acceder al primer y último elemento.\markdownRendererInterblockSeparator
{}\markdownRendererInputFencedCode{./_markdown_main/b0604721c19cea4517abe1909ed46605.verbatim}{cpp}{cpp}\markdownRendererInterblockSeparator
{}
\markdownRendererSectionEnd 
\markdownRendererSectionEnd \markdownRendererSectionBegin
\markdownRendererHeadingTwo{Estructuras Set}\markdownRendererInterblockSeparator
{}\markdownRendererSectionBegin
\markdownRendererHeadingThree{Sets y multisets}\markdownRendererInterblockSeparator
{}\markdownRendererSectionBegin
\markdownRendererHeadingFour{Set}\markdownRendererInterblockSeparator
{}Se basa en un árbol de búsqueda binario equilibrado y sus operaciones funcionan en $O(log n)$ tiempo.\markdownRendererInterblockSeparator
{}
\markdownRendererSectionEnd \markdownRendererSectionBegin
\markdownRendererHeadingFour{Unordered set}\markdownRendererInterblockSeparator
{}Se basa en una tabla hash y sus operaciones funcionan, por término medio, en $O(1)$.\markdownRendererParagraphSeparator
{}Puesto que son utilizados de la misma manera, nos centramos en la estructura del set en los siguientes ejemplos.\markdownRendererParagraphSeparator
{}La función \markdownRendererEmphasis{insert} añade un elemento al set, la función \markdownRendererEmphasis{count} devuelve el número de ocurrencias de un elemento en el set, y la función \markdownRendererEmphasis{erase} elimina un elemento del set.\markdownRendererInterblockSeparator
{}\markdownRendererInputFencedCode{./_markdown_main/86be2ff716606b37d8feba4f2822b18c.verbatim}{cpp}{cpp}\markdownRendererInterblockSeparator
{}Una propiedad importante de los sets es que todos sus elementos son distintos, por lo que la función \markdownRendererCodeSpan{ìnsert} nunca agregará un elemento si ya se encuentra dentro del set.\markdownRendererParagraphSeparator
{}La función \markdownRendererCodeSpan{count} siempre devuelve 0 (si el elemento no está en el conjunto) o 1 (si está presente)\markdownRendererInterblockSeparator
{}\markdownRendererInputFencedCode{./_markdown_main/dd6d5e310d8552bb47b25f862d6dd33d.verbatim}{cpp}{cpp}\markdownRendererInterblockSeparator
{}\markdownRendererInputFencedCode{./_markdown_main/d3fea54ea40b9b3fc5651cafecf9d57d.verbatim}{cpp}{cpp}\markdownRendererInterblockSeparator
{}La función find(x) devuelve un iterador que apunta al elemento cuyo valor es x. Si el conjunto no contiene x, el iterador devuelto será end()\markdownRendererInterblockSeparator
{}\markdownRendererInputFencedCode{./_markdown_main/771f029ab20dfa9f02b3d84340269b7a.verbatim}{cpp}{cpp}\markdownRendererInterblockSeparator
{}$\clearpage$\markdownRendererInterblockSeparator
{}
\markdownRendererSectionEnd \markdownRendererSectionBegin
\markdownRendererHeadingFour{Sets ordenados}\markdownRendererInterblockSeparator
{}La principal diferencia entre las dos estructuras es que set es ordenado, mientras que unordered_set no lo es.\markdownRendererInterblockSeparator
{}\markdownRendererInputFencedCode{./_markdown_main/ea4017e4bb893ff88c57b5e99398d2cf.verbatim}{cpp}{cpp}\markdownRendererInterblockSeparator
{}\markdownRendererInputFencedCode{./_markdown_main/4e17460a6b55baeccf842553c393b26a.verbatim}{cpp}{cpp}\markdownRendererInterblockSeparator
{}
\markdownRendererSectionEnd \markdownRendererSectionBegin
\markdownRendererHeadingFour{Multisets}\markdownRendererInterblockSeparator
{}Un multiset es un conjunto que puede tener varias copias del mismo valor.\markdownRendererSoftLineBreak
{}C++ incluye las estructuras multiset y unordered_multiset, que son análogas a set y unordered_set.\markdownRendererParagraphSeparator
{}El siguiente código agrega tres copias del valor 5 a un multiset\markdownRendererInterblockSeparator
{}\markdownRendererInputFencedCode{./_markdown_main/267faa30facfd532071b6b7daab9c4dc.verbatim}{cpp}{cpp}\markdownRendererInterblockSeparator
{}La función erase elimina todas las copias de un valor en un multiset\markdownRendererInterblockSeparator
{}\markdownRendererInputFencedCode{./_markdown_main/40ee3d1d8a8357130a7db826760d1200.verbatim}{cpp}{cpp}\markdownRendererInterblockSeparator
{}Para eliminar una sola copia de un valor, se puede hacer de la siguiente manera\markdownRendererInterblockSeparator
{}\markdownRendererInputFencedCode{./_markdown_main/3398576c55d1e73d6922c61b3d13738a.verbatim}{cppp}{cppp}\markdownRendererParagraphSeparator
{}
\markdownRendererSectionEnd 
\markdownRendererSectionEnd \markdownRendererSectionBegin
\markdownRendererHeadingThree{Maps}\markdownRendererInterblockSeparator
{}Un map es un conjunto que consiste en pares clave–valor (key–value pairs). Las claves pueden ser de cualquier tipo de dato y no tienen que ser consecutivas.\markdownRendererParagraphSeparator
{}La biblioteca estándar de C++ contiene dos estructuras de map:\markdownRendererInterblockSeparator
{}\markdownRendererUlBeginTight
\markdownRendererUlItem \markdownRendererCodeSpan{map} se basa en un árbol binario de búsqueda balanceado, y el acceso a los elementos toma tiempo O(log n).\markdownRendererUlItemEnd 
\markdownRendererUlItem \markdownRendererCodeSpan{unordered\markdownRendererUnderscore{}map} utiliza hashing, y el acceso a los elementos toma tiempo O(1) en promedio\markdownRendererUlItemEnd 
\markdownRendererUlEndTight \markdownRendererInterblockSeparator
{}\markdownRendererInputFencedCode{./_markdown_main/0ffaee6021eeb92488ff469797ca78f3.verbatim}{cpp}{cpp}\markdownRendererInterblockSeparator
{}El siguiente código recorre e imprime todos los pares clave–valor almacenados:\markdownRendererInterblockSeparator
{}\markdownRendererInputFencedCode{./_markdown_main/97c94b60234c7585af5bdd7a38c76976.verbatim}{cpp}{cpp}\markdownRendererInterblockSeparator
{}Los elementos de un map están ordenados por sus claves, mientras que los de un unordered_map no lo están.\markdownRendererParagraphSeparator
{}Eliminar una clave de un mapa\markdownRendererInterblockSeparator
{}\markdownRendererInputFencedCode{./_markdown_main/f1f05d21f2942442e7bbad2fe5dd1b54.verbatim}{cpp}{cpp}\markdownRendererInterblockSeparator
{}Comprobar si una clave existe utilizando la función \markdownRendererCodeSpan{count}, que devuelve 1 si la clave está presente, o 0 en caso contrario.\markdownRendererInterblockSeparator
{}\markdownRendererInputFencedCode{./_markdown_main/5448f1bbd01e7c325efc93c82fa4a139.verbatim}{cpp}{cpp}\markdownRendererInterblockSeparator
{}
\markdownRendererSectionEnd \markdownRendererSectionBegin
\markdownRendererHeadingThree{Priority Queues (Colas de Prioridad)}\markdownRendererInterblockSeparator
{}Una cola de prioridad es una estructura similar a un multiset, que permite insertar elementos y obtener/eliminar rápidamente el mínimo o máximo elemento, dependiendo del tipo de cola.\markdownRendererParagraphSeparator
{}Por defecto, está implementada como un heap binario máximo, lo que permite obtener y eliminar el mayor elemento en la estructura.\markdownRendererInterblockSeparator
{}\markdownRendererInputFencedCode{./_markdown_main/24faa1d989d6861ddeac9e40fe3c76b4.verbatim}{cpp}{cpp}\markdownRendererInterblockSeparator
{}Si se desea que la cola devuelva el menor elemento (un min-heap), se puede usar un comparador:\markdownRendererInterblockSeparator
{}\markdownRendererInputFencedCode{./_markdown_main/371bf8a9bafa668161ff55a873872282.verbatim}{cpp}{cpp}\markdownRendererInterblockSeparator
{}$\clearpage$\markdownRendererInterblockSeparator
{}
\markdownRendererSectionEnd \markdownRendererSectionBegin
\markdownRendererHeadingThree{Policy-Based Sets (Conjuntos Basados en Políticas)}\markdownRendererInterblockSeparator
{}El compilador g++ ofrece estructuras adicionales que no forman parte de la biblioteca estándar de C++, conocidas como policy-based data structures.\markdownRendererSoftLineBreak
{}Para utilizarlas se deben incluir las siguientes líneas:\markdownRendererInterblockSeparator
{}\markdownRendererInputFencedCode{./_markdown_main/5989c352862dccc2e082c316fa5e46bd.verbatim}{cpp}{cpp}\markdownRendererInterblockSeparator
{}Definición para valores int:\markdownRendererInterblockSeparator
{}\markdownRendererInputFencedCode{./_markdown_main/2eeacf0ea7c4b10c70099b0538dd7e9b.verbatim}{cpp}{cpp}\markdownRendererInterblockSeparator
{}Ejemplo de uso\markdownRendererInterblockSeparator
{}\markdownRendererInputFencedCode{./_markdown_main/b0fc25f0aac9258b581a63f30096ad01.verbatim}{cpp}{cpp}\markdownRendererInterblockSeparator
{}Su ventaja es que permite acceder a los elementos por su posición ordenada.\markdownRendererInterblockSeparator
{}\markdownRendererUlBeginTight
\markdownRendererUlItem \markdownRendererCodeSpan{find\markdownRendererUnderscore{}by\markdownRendererUnderscore{}order(k)} devuelve un iterador al elemento en la posición k (0-indexado).\markdownRendererUlItemEnd 
\markdownRendererUlItem \markdownRendererCodeSpan{order\markdownRendererUnderscore{}of\markdownRendererUnderscore{}key(x)} devuelve cuántos elementos son menores que x.\markdownRendererUlItemEnd 
\markdownRendererUlEndTight \markdownRendererInterblockSeparator
{}\markdownRendererInputFencedCode{./_markdown_main/0341051fba2bf45c114efc2014fae422.verbatim}{cpp}{cpp}
\markdownRendererSectionEnd 
\markdownRendererSectionEnd 
\markdownRendererSectionEnd \markdownRendererDocumentEnd